\chapter{Reproducibility Challenges}
\label{ch:reproducibility}

Having established the landscape of existing WoZ platforms and their limitations, I now examine the factors that make WoZ experiments difficult to reproduce consistently and how software infrastructure can address them. This chapter analyzes the sources of variability identified in the WoZ literature and examines how current practices in infrastructure and reporting contribute to \emph{the Reproducibility Problem}. Understanding these challenges is essential for designing a system that supports reproducible, rigorous experimentation.

\section{Sources of Variability}

\emph{The Reproducibility Problem}, as introduced in Chapter~\ref{ch:intro}, encompasses two related challenges. The first concerns \emph{execution consistency}: whether a wizard reliably follows the same experimental script across multiple trials with different participants, producing comparable robot behavior in each. The second concerns \emph{cross-platform reproducibility}: whether the same experiment can be transferred to a different robot platform with minimal change to the implementing program. Both stem from gaps in current WoZ infrastructure and are examined in this chapter. A third interpretation of the term — independent replication of a published study by researchers at other institutions — is distinct from both and is not what this thesis evaluates. It is also worth noting that execution consistency, as defined here, corresponds to what the measurement literature sometimes calls \emph{repeatability}: the degree to which the same procedure produces consistent results when repeated across multiple trials of the same study.

In WoZ-based HRI studies, multiple sources of variability can compromise execution consistency. The wizard is simultaneously the strength and weakness of the WoZ paradigm. While human control enables sophisticated, adaptive interactions, it also introduces inconsistency. Consider a wizard conducting multiple trials of the same experiment with different participants.
 Even with a detailed script, the wizard may vary in timing, with delays between a participant's action and the robot's response fluctuating based on the wizard's attention, fatigue, or interpretation of when to act. When a script allows for choices, different wizards may make different selections, or the same wizard may act differently across trials. Furthermore, a wizard may accidentally skip steps, trigger actions in the wrong order, or misinterpret experimental protocols.

Riek's systematic review \cite{Riek2012} found that very few published studies reported measuring wizard error rates or providing standardized wizard training. Without such measures, it becomes impossible to determine whether experimental results reflect the intended interaction design or inadvertent variations in wizard behavior.

Beyond wizard behavior, the custom nature of many WoZ control systems introduces technical variability. When each research group builds custom software for each study, several problems arise. Custom interfaces may have undocumented capabilities, hidden features, default behaviors, or timing characteristics researchers never formally describe. Software tightly coupled to specific robot models or operating system versions may become unusable when hardware or software is upgraded or replaced. Each system logs data differently, with different file formats, different levels of granularity, and different choices about what to record. This fragmentation undermines both execution consistency and reproducibility. Rebuilding custom infrastructure for each study makes it nearly impossible to guarantee that wizard behavior is controlled the same way across trials. More broadly, reproducing the same experiment on a different robot platform typically requires reverse-engineering or rebuilding the original software from scratch.

Even when researchers intend for their work to be reproducible, practical constraints on publication length lead to incomplete documentation. Papers often omit exact timing parameters. Authors leave decision rules for wizard actions unspecified and fail to report details of the wizard interface. Specifications of data collection, including which sensor streams were recorded and at what sampling rate, frequently go missing. Without this information, other researchers cannot faithfully recreate the experimental conditions, limiting both direct replication and conceptual extensions of prior work.

\section{Infrastructure Requirements for Enhanced Reproducibility}

Based on this analysis, I identify specific ways that software infrastructure can mitigate reproducibility challenges:

\begin{enumerate}
\item \textbf{Guided wizard execution.} Rather than merely providing tools for wizard control, an ideal WoZ platform should actively guide wizards through scripted procedures. This means presenting actions in a prescribed sequence to prevent out-of-order execution, highlighting the current step in the protocol, recording any deviations from the script as explicit events in the data log, and supporting repeatable decision logic through clearly defined conditional branches. By constraining wizard behavior within the bounds of the experimental design, the system reduces unintended variability across trials and participants.

\item \textbf{Comprehensive automatic logging.} Manual data collection is error-prone and often incomplete. The platform should automatically record every action triggered by the wizard with precise timestamps, all robot sensor data and state changes, and timing information indicating when actions were requested, when they began executing, and when they completed. The full experimental protocol should be embedded in each log file so that researchers can recover the exact script used for any session. Note that recording precise timestamps does not imply that trials must have identical timing, since human-robot interactions naturally vary in duration; rather, the system captures what actually occurred for later analysis.

\item \textbf{Self-documenting protocol specifications.} The protocol specification itself should serve as documentation. When interaction protocols are defined using structured formats such as visual flowcharts or declarative scripts rather than imperative code, they become simultaneously executable and human-readable. Researchers can then share complete, unambiguous descriptions of their experimental procedures alongside their results.

\item \textbf{Platform-independent abstractions.} To maximize the lifespan and transferability of experimental designs, the platform must separate the high-level control logic, the sequence of wizard and robot actions, from the low-level details of how specific robots execute those behaviors. This abstraction allows experiments designed for one robot to be more easily adapted to another, extending the reproducibility of interaction designs even when the original hardware becomes obsolete.
\end{enumerate}



\section{Connecting Reproducibility Challenges to Infrastructure Requirements}

The reproducibility challenges identified above directly motivate the infrastructure requirements (R1--R6) established in Chapter~\ref{ch:background}. Inconsistent wizard behavior creates the need for real-time control mechanisms (R3) that guide wizards step by step, and for automatic logging (R4) that captures any deviations that occur. Timing errors further motivate responsive, fine-grained real-time control (R3): a wizard working with a sluggish interface introduces latency that disrupts the interaction and confounds timing analysis. Technical fragmentation forces each lab to rebuild infrastructure as hardware changes, violating platform agnosticism (R5). Incomplete documentation reflects the need for self-documenting, code-free protocol specifications (R2) that are simultaneously executable and shareable, integrated into a single workflow (R1) so that the specification and the execution environment are never separated. Finally, the isolation of individual research groups motivates collaborative support (R6): allowing multiple team members to observe and review trials enables the shared scrutiny that reproducibility requires. As Chapter~\ref{ch:background} demonstrated, no existing platform simultaneously satisfies all six requirements. Addressing this gap requires rethinking how WoZ infrastructure is designed, prioritizing reproducibility and methodological rigor as first-class design goals rather than afterthoughts.

\section{Chapter Summary}

This chapter has analyzed the reproducibility challenges inherent in WoZ-based HRI research, identifying three primary sources of variability: inconsistent wizard behavior, fragmented technical infrastructure, and incomplete documentation. Rather than treating these challenges as inherent to the WoZ paradigm, I showed how each stems from gaps in current infrastructure. Software design can systematically mitigate them through enforced experimental protocols, comprehensive automatic logging, self-documenting experiment designs, and platform-independent abstractions. These design goals directly address the six infrastructure requirements identified in Chapter~\ref{ch:background}. The following chapters describe the design, implementation, and pilot validation of a system that prioritizes reproducibility as a foundational design principle from inception.
